% Part: proof-theory
% Chapter: normalization
% Section: normalization-thm

\documentclass[../../../include/open-logic-section]{subfiles}

\begin{document}

\olfileid{pt}{nor}{thm}

\olsection{The Normalization Theorem}

!!^a{proof} is in \emph{normal} form if it has no cut segments.
Clearly this is the case iff we can apply neither a permutation
conversion nor a reduction conversion to it. To \emph{normalize}
!!a{proof} means to transform it into a normal proof by applying
permutation and reduction conversions until no cuts remain. That this
is always possible is the \emph{normalization theorem}.

\begin{thm}\ollabel{thm:normalization} Any \Log{N1i}-!!{proof}
$\delta$ normalizes, i.e., can be transformed into
!!a{proof}~$\delta'$ without cuts.
\end{thm}

\begin{proof}
  By induction on the cut rank and length of~$\delta$. The induction
  hypothesis is that any !!{proof} which either has lower cut rank, or
  the same cut rank but lower cut length, normalizes.

  If the cut length is~$0$, there are no cuts, $\delta$ is already
  normal. So assume that cut rank and cut length of~$\delta$ are~$>
  0$. Pick a topmost, rightmost maximal cut segment. If the cut
  segment has length $>1$, apply a permutation conversion. If the cut
  segment has length~$1$, apply the corresponding reduction
  conversion. This results in a new !!{proof} which either has the
  same cut rank but a lower cut length, or a lower cut rank, and the
  inductive hypothesis applies.
\end{proof}

The strategy used in the above proof (always deal with a
topmost rightmost maximal cut) suggests a specific order in which
conversions must be applied to !!a{proof} in order to result in a
normal !!{proof}. Surprisingly enough, the order actually does not
matter. Applying permutation and reduction conversions in any order
eventually results in a normal proof.

The definition of segment and cut can easily be formulated for
\Log{N2i}, and the permutation and reduction conversions for \Log{N1i}
directly translate to \Log{N2i}. Consequently, the normalization
theorem also holds for \Log{N2i}.

\begin{thm}\ollabel{thm:normalization-N2i} Any \Log{N2i}-!!{proof}
$\delta$ normalizes, i.e., can be transformed into
!!a{proof}~$\delta'$ without cuts.
\end{thm}

\end{document}
